\documentclass[11pt]{article}
\usepackage{amsmath,amssymb}
\title{General magnetic currents on the original mixed-boundary box}\author{}\date{}
\begin{document}\maketitle
\paragraph{Claim and scope.}
Let $\Omega=\prod_{i=0}^3\{0,\ldots,n_i-1\}$ be one of the actual
scalar component or ghost grids of a relative cube of side $R$.
Thus $n_i\asymp R$, the component has one Neumann and three Dirichlet
directions, and the ghost has four Dirichlet directions. Let $a$
be ANY real antisymmetric connection on its internal nearest-neighbor
edges. The magnetic Laplacian $L_\epsilon$ has hopping phases
$e^{i\epsilon a}$, the original diagonal and original killing rates.
Then the proposed bound is
\[
 0\le(\log\det L_\epsilon)''_0
                   \le C\log(2+R)\sum_{\square\subset\Omega}|da|^2 .
                                                               \tag{1}
\]
No harmonicity, own-coordinate independence, or imposed boundary
value of $a$ is assumed. Tiny grids with singleton factors are included
by omission of their absent edges. The proof below uses an auxiliary
gauge and a positive path domination, not replacement of the original
determinant or its Dirichlet boundary.

\paragraph{1. An exact auxiliary Coulomb gauge and reflection extension.}
Minimize $\|a-d\chi\|_{\ell^2(E)}^2$ over all vertex functions on the
INTERNAL graph, modulo constants. Put $b=a-d\chi$. Then
$d^*b=0$ at EVERY graph vertex, including its boundary.
This is merely a scalar unitary gauge of $L_\epsilon$; original
killing diagonals and every closed-loop current remain unchanged.
In particular $db=da$.

Let $F_i:\mathbb Z\to\{0,\ldots,n_i-1\}$ be the period-$2n_i$ fold
\[
 F_i(z)=z\quad(0\le z<n_i),\qquad
 F_i(z)=2n_i-1-z\quad(n_i\le z<2n_i).
\]
Pull back $b$ to oriented edges of the periodic lattice:
$B(z,z+\hat i)$ is $b(Fz,F(z+\hat i))$ when these endpoints differ,
and zero when they coincide. The latter edges are the genuine
self-loop steps of the folded walk. $B_i$ is odd in its own edge
coordinate and even in the transverse vertex coordinates.
Literal incidence gives
\[
 d^*B(z)=d^*b(Fz)=0,\qquad
 \|dB\|_{\rm per}^2=16\|db\|^2.
\]
The periodic discrete Fourier identity therefore gives
\[
 \sum_{i,j}\|\delta_jB_i\|_{\rm per}^2
       =\|dB\|_{\rm per}^2+\|d^*B\|_{\rm per}^2
       =16\|da\|^2.                                         \tag{2}
\]
This is a finite exact identity. It keeps the normal zero edges
at reflection planes, rather than discarding their boundary energy.

\paragraph{2. Centered open currents on the folded walk.}
Let $Z_t$ be the ordinary continuous-time nearest-neighbor walk on
$\mathbb Z^4$, rate one for each oriented step, started at $x\in\Omega$,
and $Y_t=FZ_t$. $Y$ is the all-reflecting graph walk.
Its $b$-current is EXACTLY the $B$-current on $Z$.
Fix $c_i(x)=B_i(x)$ and let
\[
 M_t=J_t-\sum_i c_i(x)(Z_t^i-x_i).
\]
By $d^*B=0$, this is a martingale. For the periodic walk, translation
and the spectral inequality $1-e^{-s\lambda}\le s\lambda$ give
\[
 \sum_{x\in\Omega}E_x M_t^2
 \le C\int_0^t(1+s)\,ds\sum_{i,j}\|\delta_jB_i\|_{\rm per}^2
 \le Ct(1+t)\|da\|^2.                                      \tag{3}
\]
Indeed the two directional coefficients in the bracket are
$B_i(Z_s)-B_i(x)$ and $B_i(Z_s-\hat i)-B_i(x)$.
Sum over the whole periodic cell to dominate the original start set.
The first difference costs $Cs\|\nabla B_i\|^2$ and the unit shifted
one costs $C(1+s)\|\nabla B_i\|^2$. No pointwise coefficient bound
or fourth moment of the connection is used.

The endpoint gauge needed for a closed loop is defined on the
ORIGINAL graph: $\phi_x(y)=\sum_i c_i(x)(y_i-x_i)$. Thus
\[
 J_t-\phi_x(Y_t)
       =M_t+\sum_i c_i(x)(Z_t^i-Y_t^i).                       \tag{4}
\]
The folding term is essential. Dropping it would silently discard
the reflected boundary drift.

\paragraph{3. Normal-edge Hardy control of that folding term.}
Write $d_i(x)=\min(x_i,n_i-1-x_i)$.
The fold is one-Lipschitz and fixes $x$, so
$|Z_t^i-Y_t^i|\le2|Z_t^i-x_i|$, and the difference can be nonzero only
after the unfolded endpoint has crossed one of the bounding planes.
The elementary symmetric-walk Chernoff estimate and
$E|Z_t^i-x_i|^4\le C(t+t^2)$ imply, for $t\ge1$,
\[
 E_x|Z_t^i-Y_t^i|^2
 \le Ct\exp[-c\min(d_i(x)^2/t,d_i(x))].                       \tag{5}
\]
For $0<t\le1$, the bound $Ct e^{-c d_i(x)}$ follows directly by
dominating the displacement by the Poisson jump count and summing
its factorial tail. These bounds remain valid for multiple folds.

For each fixed transverse coordinate, $B_i$ has a ZERO normal edge
at either bounding plane. Telescoping from the nearer such edge gives
\[
 |c_i(x)|^2\le C(d_i(x)+1)
       \sum_{\text{normal segment from that plane to }x}
                                      |\delta_i B_i|^2 .
\]
After multiplication by (5), reverse the segment sum. The elementary
one-dimensional shell sums are
\[
 \sum_{d\ge0}(d+1)e^{-c\min(d^2/t,d)}\le C(1+t)\quad(t\ge1),
 \qquad
 \sum_{d\ge0}(d+1)e^{-cd}<\infty .
\]
Consequently
\[
 \sum_{x\in\Omega}\sum_i |c_i(x)|^2
              E_x|Z_t^i-Y_t^i|^2
             \le Ct(1+t)\|da\|^2.                           \tag{6}
\]
This pays the normal boundary correction with the SAME original
curl energy through (2), uniformly in every boundary value.

Combining (3)--(6) yields
\[
 \sum_x E_x|J_t-\phi_x(Y_t)|^2
                         \le Ct(1+t)\|da\|^2.                \tag{7}
\]

\paragraph{4. Return to the original killed/reflected loops.}
The original process is the all-reflecting graph walk weighted by
$\exp[-\int_0^t V_{\rm kill}(Y_s)ds]$, where $V_{\rm kill}\ge0$
is the ORIGINAL missing-edge killing potential.
Thus (7) also dominates its surviving centered-current moment.
No altered boundary current on a killing jump is needed here:
only surviving paths contribute to the loop.

Split a closed original path of length $t$ into two halves and reverse
the second half. Both now start at $x$ and end at the same $y$.
The same $\phi_x(y)$ cancels between their currents.
The original positive return kernel satisfies
\[
 \sup_{x,y}p_{t/2}(x,y)\le
             C(1+t)^{-2}e^{-ct/R^2}.
\]
The exact half-loop argument and (7) therefore show
\[
 {\cal A}(t):=\sum_x E_x[J_t^2;Y_t=x,\ {\rm survival}]
 \le Ct(1+t)(1+t)^{-2}e^{-ct/R^2}\|da\|^2.                    \tag{8}
\]
The magnetic Feynman--Kac expansion gives
$-\partial_\epsilon^2\operatorname{Tr}e^{-tL_\epsilon}|_0
={\cal A}(t)$, and hence
\[
 (\log\det L_\epsilon)''_0=\int_0^\infty {\cal A}(t)\frac{dt}{t}.
\]
Bound (8) is integrable at zero and yields (1).
The proof has kept the actual return kernel and its Dirichlet gap;
the auxiliary reflecting walk was used only for the nonnegative
open-moment estimate.

\paragraph{5. Proposed leading full-harmonic implication.}
For an actual Cartan Maxwell-harmonic boundary tangent $h$, the
same-component scalar phases of the full Wilson charged Hessian,
after its exact scalar gradient gauge, are
$a^\mu_\nu(x)=h_\nu(x)+h_\nu(x+\hat\mu)$.
Their literal plaquette curvatures are
$(dh)_{\nu\rho}(x)+(dh)_{\nu\rho}(x+\hat\mu)$.
The ghost connection is $2h$. Bounded incidence therefore gives
$\sum_\mu\|da^\mu\|^2+\|d(2h)\|^2\le C E(h)$,
including the available original boundary plaquettes.

Applying (1) to these scalar operators, and the
local spin/neutral/diagonal estimates of the harmonic shear leading kinetic bound calculation, would give
\[
 |F_{\rm quad}''[h,h]|=|h^*\widehat K h|
                 \le C\log(2+R)E(h)                         \tag{9}
\]
for GENERAL actual harmonic boundary data, including the explicit
outside-shear field of the mixed harmonic shear gauge quotient calculation. Canonical source units contribute
$B^{-1}$. All determinant multiplicities and raw source-zero scalars
remain retained. Global color symmetry then supplies all three
infinitesimal colors.

The unfolding, the normal-edge boundary term and the full scalar-phase
interface of this extension still require independent verification.
Even if (9) holds, it is an UNCUT leading kinetic statement. It does
not estimate the complete two-loop coefficient $C_R''$, the matched
cutoff packet, or the original nonlinear remainder, and does not
close S94.
\paragraph{Finite checks (not substitutes for the bounds).}
The finite algebraic calculation verifies exact boundary
co-closedness, the factor-sixteen Gaffney identity, and three hundred
unfolded/folded current endpoint identities on three original finite
rectangles. The finite algebraic calculation
checks the actual full charged Wilson incidence on $R=2,3$ for an
outside-shear harmonic cubic: the original force and the general
scalar-phase residual both vanish exactly. These test the finite
interfaces; the uniform Hardy and moment estimates are proved above.
\end{document}
